\documentclass[10pt]{article}
\usepackage[margin=1.9cm]{geometry}
\usepackage{booktabs}
\usepackage{array}
\usepackage{float}
\usepackage{url}
\usepackage{graphicx}
\usepackage[colorlinks=true, urlcolor=blue]{hyperref}
\usepackage{xeCJK}
\setCJKmainfont{Droid Sans Fallback}
\setlength{\floatsep}{6pt}
\setlength{\textfloatsep}{8pt}

\title{License Plate Detection and Recognition on CCPD2019\\ \large Technical summary}
\date{10 November 2026}

\begin{document}
\maketitle
\vspace*{-1.4cm}

\section{Problem and question}

This project develops a pipeline that detects and recognizes Chinese
license plates in CCPD2019 images. We evaluate two controlled
improvements over an off-the-shelf baseline --- fine-tuning a
pretrained YOLO11s detector on our 5{,}000-image subset, and replacing
the general-purpose EasyOCR engine with a task-specific convolutional
recurrent neural network (CRNN) --- and ask whether they improve
end-to-end plate recognition.

\section{Dataset}

We built a reproducible subset of CCPD2019: 5{,}000 source images
stratified by plate height and the two tilt values encoded in the
filenames. Each source image was generated at
three scales (1.0, 0.75, and 0.5), producing 15{,}000 image variants.
The subset was split into 3{,}500 training, 700 validation, and 800 test
images using seed 42, before resizing, so that all variants of one image
remain in the same split and cross-scale data leakage is prevented.

Plate size categories used training-set height thresholds: small
($\le 83$~px), medium (83--112~px), large ($> 112$~px). We converted the
CCPD filename annotations into normalized YOLO boxes and plate crops
with text labels. Three OCR preprocessing modes support the ablation:
none, grayscale + CLAHE, and Canny. Primary experiments use
original-scale images; downscaled variants support the scale-robustness
evaluation.

\section{Method}

\textbf{Detection.} We used YOLO11s to locate license plates before
recognition. The detector predicts a single class (license plate); the
highest-confidence box is cropped and passed to the recognizer. As a
baseline, we used the frozen
pretrained detector from MorseTechLab
(\href{https://huggingface.co/morsetechlab/yolov11-license-plate-detection}{the Hugging Face repository},
\texttt{license-plate-finetune-v1s.pt}). We fine-tuned it on 3{,}500
original-scale training images (640~px, batch 16, AdamW, learning rate
0.002) and selected the best checkpoint by validation mAP@0.5 (epoch 4
of 24).

\textbf{Recognition.} As a baseline, we used the frozen general-purpose
EasyOCR engine. As a second condition, we trained a plate-specific CRNN
(CNN + bidirectional LSTM + CTC, 2{,}999{,}419 parameters) from scratch:
each crop is converted to grayscale and resized onto a 48~$\times$~256
canvas while preserving its aspect ratio, and the classifier predicts the
58 training-set characters plus the CTC blank. Training used 3{,}500
original-scale crops without augmentation (50 epochs, Adam, learning rate
0.001, batch 32, seed 42); the epoch-39 checkpoint was selected by
validation character accuracy. Preprocessing mode and confidence
threshold (0.50) were selected on validation: no preprocessing.

\section{Results}

Table~\ref{tab:main} summarizes the main test metrics (800 images and
800 crops).

\begin{table}[H]
\centering
\small
\begin{tabular}{@{}lrrrr@{}}
\toprule
Configuration & mAP@0.5 & Detection F1 & Character accuracy & Exact plates \\
\midrule
Frozen YOLO11s + EasyOCR & 0.9073 & 0.9167 & 0.3240 & 0 / 800 \\
CRNN, ground-truth crops & --- & --- & 0.7593 & 267 / 800 \\
CRNN + frozen YOLO11s & 0.9073 & 0.9167 & 0.6710 & 148 / 800 \\
CRNN + fine-tuned YOLO11s & 0.9949 & 0.9714 & 0.7661 & 296 / 800 \\
\bottomrule
\end{tabular}
\caption{Test metrics of the evaluated configurations ($n = 800$).}
\label{tab:main}
\end{table}

Fine-tuning raised detection recall from 0.8938 to 0.9975 and precision
from 0.9408 to 0.9466. On the ground-truth crops, CRNN outperformed
EasyOCR with 0.7593 versus 0.3474 character accuracy and 267 versus 0
exact plates.

\textbf{Preprocessing ablation.} Table~\ref{tab:ablation} reports the
validation character accuracy of the three preprocessing modes; the
selected policy used no preprocessing and a confidence threshold of
0.50 in all protocol groups.

\begin{table}[H]
\centering
\small
\begin{tabular}{@{}lr@{}}
\toprule
Preprocessing mode & Character accuracy \\
\midrule
No preprocessing & 0.7558 \\
Grayscale + CLAHE contrast & 0.7469 \\
Canny edge detection & 0.4575 \\
\bottomrule
\end{tabular}
\caption{Preprocessing ablation (validation character accuracy, CRNN).}
\label{tab:ablation}
\end{table}

\textbf{Scale robustness.} The 0.75 and 0.5 scale variants of the test
set were evaluated with the same weights and the validation-selected
policy, without retraining (Table~\ref{tab:scale}).

\begin{table}[H]
\centering
\small
\setlength{\tabcolsep}{6pt}
\begin{tabular}{@{}lrrr@{}}
\toprule
Metric & 1.0 & 0.75 & 0.5 \\
\midrule
Detection mAP@0.5 / F1, frozen & 0.9073 / 0.9167 & 0.9218 / 0.9271 & 0.9294 / 0.9298 \\
Detection mAP@0.5 / F1, fine-tuned & 0.9949 / 0.9714 & 0.9949 / 0.9726 & 0.9949 / 0.9708 \\
\midrule
CRNN, ground-truth crops & 0.7593 (267) & 0.7613 (265) & 0.7610 (268) \\
CRNN + frozen YOLO11s & 0.6710 (148) & 0.6713 (157) & 0.6721 (144) \\
CRNN + fine-tuned YOLO11s & 0.7661 (296) & 0.7629 (297) & 0.7555 (279) \\
\bottomrule
\end{tabular}
\caption{Test metrics by scale (recognition: character accuracy; exact
plates in parentheses).}
\label{tab:scale}
\end{table}

Both stages remain flat across scales: the fine-tuned detector keeps
mAP@0.5 of 0.9949 at all three scales, and CRNN accuracy on the
ground-truth crops varies within 0.7593--0.7613 because the recognizer
resizes every crop to the fixed 48~$\times$~256 canvas.

\section{Interpretation}

The end-to-end gap of the frozen detector to the OCR-only performance
(0.6710 versus 0.7593 character accuracy) closes almost entirely after
fine-tuning (0.7661): missed detections are the dominant source of
end-to-end error. The task-specific CRNN beats the general-purpose
EasyOCR, which recognized no plate exactly, and because every crop is
resized to the fixed 48~$\times$~256 canvas, recognition accuracy is set
by crop content, not absolute image scale.

\section{Failure analysis}

We inspected representative failures from the test set and the
preprocessing ablation (Table~\ref{tab:failures}).

\begin{table}[H]
\centering
\small
\setlength{\tabcolsep}{4pt}
\setlength{\tabcolsep}{5pt}
\begin{tabular}{@{}p{4.4cm}p{5.2cm}p{5.4cm}@{}}
\toprule
Stage (sample) & Observation & Likely cause \\
\midrule
Frozen detector (\texttt{6f3c772c\ldots}) & Large plate, 19$^{\circ}$
tilt, no prediction at any score & Weak prior for large plate aspect
ratios in the pretrained model, especially at 10--19$^{\circ}$ tilt \\
Fine-tuned detector (\texttt{356b4ca6\ldots}) & 11 boxes: one true (0.84)
and ten false positives & Fine-tuning removed the misses but left
low-confidence false positives on cluttered scenes \\
CRNN recognition (\texttt{588e616d\ldots}) & Ground truth
\texttt{苏E050R8}, prediction \texttt{05oRo菲} (large, 19$^{\circ}$ tilt) &
Fixed 48~$\times$~256 input and training dominance of small, near-0$^{\circ}$
plates leave large tilted glyphs out of distribution \\
CRNN + Canny (\texttt{ce92ecf7\ldots}) & Ground truth \texttt{浙JD389L},
prediction \texttt{皖AU0367} (small, 3$^{\circ}$ tilt) & The Canny edge map
leaves only hollow stroke outlines; the lost fill and shading make the
glyphs unreadable \\
\bottomrule
\end{tabular}
\caption{Representative failure cases and their likely causes.}
\label{tab:failures}
\end{table}

The failures concentrate on large plates and the 10--19$^{\circ}$ tilt
bin, where the OCR-only character accuracy drops to 0.5658 and 0.5193.
Larger inputs and size/tilt-balanced training are the natural next
steps.

\section{Reproducibility}

Code: \href{https://github.com/S1norin/inno-detecting-license-plates}{the project repository}.

Demo video: \href{https://drive.google.com/file/d/1KBiiD-dvhvzsI7yWwDc7KqwM6jp585ez/view}{demo video (Google Drive)}.

Weights: \href{https://huggingface.co/Sinorin/inno-detecting-license-plates-weights}{the weights repository}
(fetched with SHA-256 verification). Python 3.12, single-threaded CPU
inference; training 4{,}154.9~s (CRNN) and 1{,}439.8~s (detector);
inference about 83--86~ms per image. Reproduction commands:
\texttt{docs/md/final\_results.md}.

\section{Members' contributions}

\begin{itemize}
\setlength{\itemsep}{1pt}
\item \textbf{Andrei Polevoi} --- Data \& preprocessing: subset, size buckets, splits, labels, preprocessing module.
\item \textbf{Maximilian Mifsud Bonici} --- Detection: YOLOv11 baseline + fine-tuning, mAP evaluation, failure cases.
\item \textbf{Evgenii Khovalyg} --- Recognition: CRNN training + OCR baseline, character/plate metrics, fine-tuning.
\item \textbf{Mikhail Krylov} --- Pipeline \& evaluation: end-to-end integration, thresholds, latency, sensitivity tables.
\end{itemize}

\end{document}
